\section{Discussion}
\label{ch:discussion}

This section answers each research question against the evidence of Section~\ref{ch:results}, sets out the mechanism the findings share, relates them to the prior work of Section~\ref{ch:related}, and states the scope and the implications.

\subsection{Answers to the Research Questions}
\label{sec:disc-rqs}

Each research question is answered below against the section of Section~\ref{ch:results} that establishes the answer. The answers are stated rather than re-argued; the evidence and its qualifications are where they were presented.

\textbf{RQ1} asked whether the input-embedding gradient of a frozen autoregressive sequence likelihood provides a useful proposal direction for discrete token revision, and if not, which alternative local quantities do. The first half is answered in the negative and the second constructively. Sampling with that gradient is indistinguishable from sampling with a norm-matched random direction, certified as equivalence at a margin fixed in advance over a thousand paired sequences (Section~\ref{sec:results-fallacy}), and the null holds across a sweep spanning proposal entropies from uniform to deterministic (Section~\ref{sec:results-sweep}). The mechanism has two parts: the first-order surrogate is uncorrelated with the true energy change over the entire admissible range of substitution distances, so the gradient is inapplicable rather than imprecise (Section~\ref{sec:results-linradius}), and the derivative is structurally blind to how well a candidate fits its own left context, exactly so at the final position where it is identically zero while the energy difference between candidate final tokens is exactly the model's conditional log-ratio (Section~\ref{sec:results-lasttoken}). The alternatives are named and measured. The relaxed token-indicator derivative restores the discarded term in closed form, correlates at $\rho = 0.60$ to $0.73$ at admissible distances, and inside the same sampler recovers $40$ percent of corrupted tokens on GPT-2 Large and $41$ on Llama-3 (Section~\ref{sec:results-onehot}); the model's own next-token conditional recovers $23.5$ percent in the exact-energy chain, and a bidirectional proposal more still (Table~\ref{tab:mlm}). Under its proposal name, the question of whether energy-based recovery improves on ordinary autoregressive decoding resolves in two directions at once: negatively for the gradient-guided samplers, affirmatively for exact-energy rescoring, whose top-$k$ rerank reaches a final KL of $4.43$ (Section~\ref{sec:results-baselines}).

\textbf{RQ2} asked about the role of the Metropolis--Hastings correction in each sampler, and it plays opposite roles, which is itself the finding. For the discrete sampler the correction is what makes the chain work at all, because its proposal is numerically uniform at every step: without it the chain is a random walk that does not converge, and with it the accept/reject filter supplies the entire selection pressure (Section~\ref{sec:results-quench}). Here the exact accept-reject step and empirical performance align, for a reason unflattering to the method, since what works is the filter and not the guidance. For the continuous sampler the same correction is disabling, rejecting almost every boundary-crossing move, and the decomposition attributes the collapse to the proposal term rather than the target term: the moves are good and their reversibility cannot be verified (Section~\ref{sec:results-mh}), because the drift derived from the differentiable pathway is not Lipschitz across a cell boundary while the projected energy is piecewise constant and never differentiated. In the implementations and landscapes tested here, omitting the correction makes the update behave as an early-stopped stochastic optimizer rather than an exact sampler from the stated target. How far that reframes any particular published system is not something this thesis measured.

\textbf{RQ3a} asked whether a GFlowNet trained on the frozen likelihood learns a policy that generates high-reward text. It does, and that is the problem: the policy attains high reward by exploiting the reward rather than by writing well. Three failure modes appear, capacity collapse, length collapse under the unnormalized reward, and reward hacking under the normalized one, and the two settings of the single length parameter produce two opposite degeneracies (Section~\ref{sec:results-gfn}). The finding covers the variants trained here and not amortized inference in general, since a GFlowNet's failure can also turn on reward design, training stability, policy parameterization, capacity, exploration, termination handling, or reward scaling.

\textbf{RQ3b} asked whether the energy induced by that tuning becomes more navigable by the local input-embedding surrogate than the untuned base. It does not. Tuning moved the energy by tens of nats and dropped the base-to-tuned correlation of sequence log-likelihoods as low as $0.77$, yet Langevin sampling on the tuned models is indistinguishable from sampling on the base, and the linearization correlation on the tuned energies stays flat (Section~\ref{sec:results-gfn}). Amortization and energy smoothing are different objectives, and this is the reason the two halves of RQ3 are reported separately: the first is a statement about a learned policy, the second about the geometry that policy leaves behind.

\textbf{The extension question E} asked whether an additive constraint can steer generation on this energy landscape. On the frozen autoregressive landscape it cannot. The trustworthy paired contrast between a real and a randomized constraint direction is essentially zero on the discrete sampler, and the one setting in which a constraint-direction signal does appear, the constraint-only arm on a continuous baseline, is entangled with the departure from the fluent manifold that a constraint-only objective causes, so a classifier scoring that text is scoring a distribution it was not trained for (Section~\ref{sec:results-constrained}). On a score-trained diffusion landscape, steering became partly possible under a fluency trust region, in one direction per setting, with the reachable direction tracking how far the guiding and judging instruments agree on the text being scored (Appendix~\ref{app:guided}). These are two answers about two landscapes and are kept apart; the second is a repair experiment, labelled exploratory where it is reported.

\subsection{One Problem, Four Mechanisms}
\label{sec:disc-unified}

It would be tidy to say that a single defect explains every failure in this thesis, and it would be wrong. The failures share a high-level problem and then diverge into mechanisms that are genuinely distinct, and the honest account is a hierarchy rather than a unification.

The shared problem is in the original plug-and-play construction, which asks the raw frozen likelihood to play three roles at once: a quality objective whose optimum is the text one wants, a locally navigable energy whose neighbourhood structure a proposal can exploit, and a reward in proportion to which a policy can be trained to sample. Teacher forcing optimizes each conditional on a correct prefix and constrains none of the three, and the experiments show that the roles come apart in different ways.

Four lower-level mechanisms are involved, and they are not manifestations of one another. The first is a \textbf{proposal-information failure}: the input-embedding derivative omits the candidate's direct conditional-fit term, which is why the surrogate carries no rank information and why supplying the missing term repairs it (Section~\ref{sec:results-linradius}). The second is a \textbf{continuous-state correctness failure}: a nearest-neighbour projection makes the drift discontinuous across cell boundaries, so the reverse-proposal density vanishes and the correction rejects exactly the moves that change a token (Section~\ref{sec:results-mh}). The third is an \textbf{objective-quality failure}: high sequence likelihood corresponds to degenerate text, so the target of the whole procedure is not the target one would want (Section~\ref{sec:results-trap}). The fourth is an \textbf{amortized-training failure}: the GFlowNet variants collapse along whichever axis the reward's implicit incentives push, which is a matter of reward shape and training dynamics and not of any gradient (Section~\ref{sec:results-gfn}). The independence is easy to check. The GFlowNet never differentiates its reward, so its collapse cannot be caused by a missing self term in a derivative it does not take; the likelihood trap concerns the location of the optimum and would persist under a perfect proposal; and the projection geometry belongs to the continuous relaxation alone, the discrete sampler never leaving the token manifold.

What runs through the whole study, and is the closest thing to a single thread, is a statement about proposals rather than mechanisms: the proposal cannot see token fitness. It cannot see it when the derivative is taken in the wrong coordinates, cannot act on it when the correction rejects every move that would use it, and cannot be given it by a reward that is itself maximized by degenerate text. The ladder of Table~\ref{tab:mlm} is the positive form of the same statement, recovery rising monotonically with what the proposal may condition on. This thesis refers to the Section~\ref{sec:results-fallacy} result by the shorthand \emph{gradient fallacy}, used only here and with its scope attached: it names a bounded empirical finding about the input-embedding derivative, not a conceptual error across the field and not a property of autoregressive likelihoods as such.

\subsection{Connections to the Related Work}
\label{sec:disc-related}

The likelihood trap measured in Section~\ref{sec:results-trap} is a direct, quantitative reproduction on the study's own models of the neural-text-degeneration phenomenon that \citet{holtzman2020curious} identified: that the most probable sequences under a neural language model are degenerate, and that maximization-based decoding produces worse text than sampling. Section~\ref{ch:related} presented that as a known hazard of decoding. The contribution here is to re-derive it as an obstacle for energy-based sampling specifically, since in that framing the degeneracy of high-probability text says that the \emph{target} of the sampling procedure is undesirable rather than that a decoding heuristic should be avoided.

The anisotropy results of Section~\ref{sec:results-aniso} confirm on GPT-2 Large and Llama-3 the representation-degeneration and embedding-anisotropy phenomena documented by \citet{gao2019representation} and \citet{ethayarajh2019contextual}, and add a consequence those works did not draw: the anisotropy scale differs by roughly a factor of three between the two models, so a single step-size schedule cannot transfer, which makes the samplers model-specific in a way that qualifies their claim to be plug-and-play.

The Metropolis--Hastings breakdown of Section~\ref{sec:results-mh} is the empirical counterpart to a gap noted in Section~\ref{sec:rel-energy}: that the energy-based decoding literature frequently omits the correction, and that where it is included, the difficulty of computing the reverse proposal on a projected landscape is seldom examined. This thesis implements the correction faithfully, following the convergence theory of \citet{roberts1996exponential}, and measures what the common omission conceals. On that reading, the correction-free update of the mechanism that COLD \citep{qin2022cold} and MuCoLa \citep{kumar2022gradient} share, and that the continuous sampler of this thesis follows, is better understood as a biased optimizer than as a sampler from the stated energy. That is a statement about the mechanism as reimplemented and measured here, not a reinterpretation of those systems' published results, which involve tasks, tuning, constraints and evaluation procedures this study did not reproduce. The same reading applies to the annealed stochastic-gradient Langevin dynamics of \citet{welling2011bayesian}. Its guarantee relies on an annealed schedule whose step size shrinks the proposal onto a local neighbourhood; Section~\ref{sec:results-quench} shows that on this landscape the schedule does not achieve that, because the proposal remains near-uniform over the vocabulary throughout, so the uncorrected chain neither localizes nor converges.

The GFlowNet results qualify, without contradicting, the picture presented by \citet{hu2024amortizing}, whose formulation and codebase this thesis builds on. Their work showed that amortized sampling from a frozen-model posterior can yield diverse, transferable samples where the reward is well-behaved, and nothing here disputes that; the narrower finding is that when the reward is the raw frozen likelihood used over free-form sequences, the learned policy collapses in the ways Section~\ref{sec:results-gfn} documents. The axis that reconciles the two is whether the reward defines a landscape with a coherent-text optimum.

The gradient-free samplers isolated in Section~\ref{sec:rel-energy}, Mix-and-Match \citep{mireshghallah2022mixmatch} and the twisted sequential Monte Carlo of \citet{lew2023sequential} and \citet{zhao2024probabilistic}, confirm rather than contradict the central result: each samples from a frozen-model energy without differentiating it, and each works, as the Gibbs baseline of Section~\ref{sec:results-baselines} does in-platform on the identical energy. The frozen likelihood is a serviceable object to evaluate and to search without gradients, and an unusable one to differentiate in its input-embedding coordinates, which is why the claim is narrowed to that derivative rather than to the energy.

\subsection{Scope and Limitations}
\label{sec:disc-scope}

The thesis shows that in the settings tested, the frozen likelihood of a standard autoregressive language model is a poor energy for gradient-guided and amortized sampling on discrete text when the gradient is taken with respect to its input embeddings, and it explains why in terms of the coordinates that derivative is taken in, the geometry of the token embedding space, and what the resulting proposal is permitted to condition on. It does not show that energy-based controllable generation is impossible in general.

Four limitations bear on how far the central result can be pushed. First, and most importantly, the claim is about one Jacobian slice: the relaxed token-indicator derivative of the same frozen likelihood carries a usable ranking signal and, in the same sampler, recovers up to $40$ percent of corrupted tokens where the input-embedding gradient reaches at most two (Section~\ref{sec:results-onehot}). The thesis therefore refutes the premise for a sampler that differentiates the input embedding of a likelihood whose target token enters as a discrete index, and not for gradient-guided discrete sampling as such. That scope is narrower than the family this work set out to test, and the difference is worth stating precisely. MuCoLa \citep{kumar2022gradient} substitutes the continuous vector into the output softmax itself, defining $P(\tilde{\bm{e}}_{n+1} \mid \tilde{\bm{e}}_{1:n})$ by replacing the looked-up target embedding with the state, so its gradient reaches the token's own score directly and not only through the following hidden states; COLD \citep{qin2022cold} carries a soft sequence of vocabulary logits whose fluency term is a soft cross-entropy against the model's reference distribution, differentiable in the same way. Neither discards the self term. The energy implemented here scores the masked position by gathering at the projected token index, and it is that construction, not theirs, that the null characterizes. The finding is correspondingly reframed rather than weakened: a natural implementation choice inside the embedding-space approach destroys the ranking signal, the published methods avoid it by relaxing the target as well as the input, and Section~\ref{sec:results-onehot} measures what that choice is worth. The convergence is close enough to be worth recording. MuCoLa's self term contributes a gradient of $\bm{h}_n$, so it ranks a candidate $v$ against the incumbent by $\bm{h}_n^\top(\bm{e}(v) - \bm{e}(x_i))$, the difference of their logits, while the relaxed token-indicator self term contributes $\log p(v \mid \bm{x}_{<i}) - \log p(x_i \mid \bm{x}_{<i})$; the two differ only by the shared normalizer, which cancels in the difference. What the continuous sampler of this thesis does reproduce faithfully is the state geometry, a continuous embedding state with a nearest-neighbour projection onto the embedding table after every update, and it is that geometry the Metropolis--Hastings result of Section~\ref{sec:results-mh} concerns. The result holds across architectures, reaching $41$ percent on Llama-3 8B, but carries two qualifications rather than an unqualified generality. It requires a configuration in which the surrogate and not the distance term shapes the proposal, and neither model's calibrated setting is such a configuration, so what is established is that the signal is present and usable, not that any off-the-shelf schedule will use it. And the cross-architecture evidence is weaker than the GPT-2 evidence: on GPT-2 the two surrogates are compared cell for cell over one step-size and temperature grid, whereas on Llama no input-embedding run was made at the surrogate-driven setting, so the comparison there is against every Llama configuration that was run rather than against a matched one. Second, the flagship ablation was measured where the proposal is numerically uniform, and it is the sweep of Section~\ref{sec:results-sweep}, at $n = 50$ per cell rather than $200$, that carries the null into regimes where the gradient is load-bearing. Third, the correction was applied asymmetrically across the compared arms in the main grid (Section~\ref{sec:meth-configs}); that cannot manufacture the null, since it penalizes the policy arm, but it did manufacture the appearance of active harm, and the corrected reruns cover three cells rather than the whole grid. Fourth, the certified equivalence rests on a thousand paired sequences drawn from independent corruptions of $282$ distinct sentences, so it is a statement about variation across corruptions and the sentence-level base remains narrow.

One formulation deserves particular care, because the thesis set out to test it. The experiments refute the premise that the frozen autoregressive sequence-likelihood gradient can supply the required local score for a Langevin sampler on a likelihood whose target token enters as a discrete index, without additional training. They do not refute the training-free premise as such. Training-free methods that do not differentiate the raw joint likelihood remained viable in this study's own results: the gradient-free Gibbs sampler of Section~\ref{sec:results-baselines} is comparable to the best gradient-guided Langevin configuration on the identical energy, and a single top-$k$ rescoring pass beats every configuration in the grid, both without any training and without any gradient. Methods that use lookahead, conditional resampling, or another representation may likewise remain viable, and Section~\ref{sec:rel-energy} reviews published members of that family.

Several further caveats bound the empirical claims. The KL-divergence metric depends on the position of the corrupted token, which is why the central results are argued from the absence of a gap between conditions rather than from small measured differences. The Llama-3 results support only qualitative cross-architecture claims, because the tokenizer and embedding scale differ from the GPT-2 family. The linearization result is that the surrogate is uniformly uninformative across the admissible range rather than informative up to a sharp radius and useless beyond it, so ``linearization radius'' summarizes a decay rather than naming a threshold. The GFlowNet result covers the variants trained here and not amortized inference in general. And the constrained-generation extension establishes the absence of reliable steering on this energy, not the impossibility of steering on a better one.

Finally, the relationship between the title of this thesis and its content. The original objective was controllable generation via faithful Langevin dynamics, in two stages: establish that energy-guided sampling on a frozen model is viable, then add the constraint term and perform sentiment and toxicity control. The first stage did not hold, and the work turned to why rather than building control on a foundation it had shown to be unsound. What it delivers is a diagnosis of the failure of a prerequisite, with the mechanisms identified, measured and cross-checked, and with the extension included as evidence of the downstream consequence, rather than a demonstration of successful control.

Several deliverables named in the proposal became moot with that turn, and each is accounted for here. Toxicity control on RealToxicityPrompts was dropped because once sentiment steering failed on the same machinery (Section~\ref{sec:results-constrained}), a second attribute on the same energy would only have repeated the negative result. The MAUVE and Self-BLEU diversity suite was dropped because those are distribution-level metrics for open-ended generation and make little sense for fixed-span infilling once the task narrowed to recovery. The planned 70B-scale judge was reduced to an 8B external judge (Section~\ref{sec:results-judge}), which was enough to rule out metric circularity without the larger model. The posterior-coverage evaluand was abandoned because the samplers never reached a regime in which coverage is a meaningful question, as the near-uniform-proposal and Metropolis--Hastings results show. And the CommonGen benchmark was not needed, ROCStories being sufficient for the diagnostic and matching the amortization reference implementation of \citet{hu2024amortizing}.

\subsection{Implications}
\label{sec:disc-future}

The most consequential implication is the one the study did not set out to find. An earlier reading of these results located the defect in the training objective and predicted that a model trained to supply a score would not fail in the same way. A score-trained diffusion model does indeed succeed where the input-embedding gradient fails, but the ladder of Table~\ref{tab:mlm} shows that the prediction was right for the wrong reason: a bidirectional masked language model, never trained with any score objective, does better still, and the frozen autoregressive model's own conditional already recovers a quarter of the tokens its gradient never reaches. What separates a working proposal from a failing one is the derivative it is built from and the context it may condition on, not the objective that trained it.

Two implications follow for practice. The first concerns what inference-time control on a frozen autoregressive model can be built from, and it is a correction of the slogan an earlier version of this work would have offered. It is not simply that one should evaluate the energy and not differentiate it, because a derivative of the same frozen likelihood does work: what one must not do is differentiate it in the input-embedding coordinates. Either differentiate the right object, taking the relaxed token-indicator derivative whose self term is one forward pass away, or bypass the derivative and propose from the model's output side directly; both routes are available on a frozen model, and Appendix~\ref{app:cost} shows that both are cheaper than the backward pass they replace. The second concerns what a further model change could buy. Since the remaining gap after the derivative is fixed is bought by right-context access, the intervention that addresses it is not a score-matching fine-tune but any procedure that gives the proposal a bidirectional view of the position it is filling, which an off-the-shelf masked language model already supplies at no training cost. Only once a landscape is known to be navigable does the additive constraint term of the plug-and-play framework become meaningful, since only then is there a fluency energy a sampler can follow for the constraint to be added to.
